
\documentclass{article}
\usepackage{colortbl}
\usepackage{makecell}
\usepackage{multirow}
\usepackage{supertabular}

\begin{document}

\newcounter{utterance}

\centering \large Interaction Transcript for game `hot\_air\_balloon', experiment `air\_balloon\_survival\_lt\_negotiation\_easy', episode 3 with mistral{-}large{-}3{-}azure.
\vspace{24pt}

{ \footnotesize  \setcounter{utterance}{1}
\setlength{\tabcolsep}{0pt}
\begin{supertabular}{c@{$\;$}|p{.15\linewidth}@{}p{.15\linewidth}p{.15\linewidth}p{.15\linewidth}p{.15\linewidth}p{.15\linewidth}}
    \# & $\;$A & \multicolumn{4}{c}{Game Master} & $\;\:$B\\
    \hline

    \theutterance \stepcounter{utterance}  
    & & \multicolumn{4}{p{0.6\linewidth}}{
        \cellcolor[rgb]{0.9,0.9,0.9}{
            \makecell[{{p{\linewidth}}}]{
                \texttt{\tiny{[P1$\langle$GM]}}
                \texttt{Dalyvaujate bendradarbiavimo derybų žaidime.} \\
\\ 
\texttt{Kartu su kitu dalyviu turite susitarti dėl vieno daiktų rinkinio, kuris bus išlaikytas. Kiekvienas iš jūsų turi savo nuomonę apie tai, kiek kiekvienas daiktas jums svarbus (svarba). Jūs nežinote, kaip kitas dalyvis vertina daiktus. Be to, jums nurodoma, kiek pastangų reikalauja kiekvienas daiktas.} \\
\texttt{Dėl rinkinio galite susitarti tik tuo atveju, jei pasirinktų daiktų bendros pastangos neviršija bendros ribos:} \\
\\ 
\texttt{LIMIT = 3712} \\
\\ 
\texttt{Čia pateikiamos atskirų daiktų pastangų vertės:} \\
\\ 
\texttt{Daikto pastangos = \{"C08": 949, "C55": 388, "A43": 79, "A67": 585, "B42": 564, "A62": 230, "B19": 580, "C43": 84, "C73": 976, "C10": 274, "B36": 374, "A18": 913, "C61": 303, "C67": 578, "C42": 548\}} \\
\\ 
\texttt{Čia pateikiama jūsų asmeninė nuomonė apie kiekvieno daikto svarbą:} \\
\\ 
\texttt{Daikto svarbos reikšmės = \{"C08": 138, "C55": 583, "A43": 868, "A67": 822, "B42": 783, "A62": 65, "B19": 262, "C43": 121, "C73": 508, "C10": 780, "B36": 461, "A18": 484, "C61": 668, "C67": 389, "C42": 808\}} \\
\\ 
\texttt{Tikslas:} \\
\\ 
\texttt{Jūsų tikslas {-} susitarti dėl bendro daiktų rinkinio, kuris būtų jums kuo naudingesnis (t. y. maksimaliai padidintų bendrą svarbą JUMS), neviršijant LIMIT. Neprivalote kiekvienoje žinutėje pateikti PASIŪLYMAS {-} galite tiesiog derėtis. Visos taktikos yra leidžiamos!} \\
\\ 
\texttt{Sąveikos protokolas:} \\
\\ 
\texttt{Žinutėje galite naudoti tik toliau nurodytus struktūrinius formatus:} \\
\\ 
\texttt{PASIŪLYMAS: \{'A', 'B', 'C', …\}} \\
\texttt{Siūlykite išlaikyti būtent šiuos daiktus.} \\
\texttt{ATSISAKAU: \{'A', 'B', 'C', …\}} \\
\texttt{Aiškiai atmeskite oponento pasiūlymą.} \\
\texttt{ARGUMENTAS: \{'...'\}} \\
\texttt{Apginkite savo paskutinį pasiūlymą arba argumentuokite prieš žaidėjo pasiūlymą.} \\
\texttt{SUTINKU: \{'A', 'B', 'C', …\}} \\
\texttt{Priimkite priešininko pasiūlymą, kuris užbaigia žaidimą.} \\
\\ 
\\ 
\texttt{Taisyklės:} \\
\\ 
\texttt{Galite SUTINKU tik su pasiūlymu, kurį kita šalis užregistravo per PASIŪLYMAS.} \\
\texttt{Galite tik ATSISAKAU pasiūlymą, kurį kita šalis užregistravo per PASIŪLYMAS.} \\
\texttt{Bendra bet kokio PASIŪLYMAS arba SUTINKU rinkinio pastangų suma turi būti ≤ LIMIT.} \\
\texttt{NESKELBKITE savo paslėptų svarbos balų.} \\
\texttt{Po struktūrinio formato žymos turi sekti dvitaškis ir tarpas. Argumentas turi būti python rinkinys, kuriame yra 0 ar daugiau eilučių.} \\
\texttt{Taigi jis turi būti tokios formos: "TAG: \{...\}"} \\
\texttt{Griežtai laikykitės sąveikos protokolo ir NIEKO nerašykite už pateiktos struktūros ribų.} \\
\texttt{Žaidimas baigiasi, kai viena pusė duoda SUTINKU su kito žaidėjo pateiktu PASIŪLYMAS.} \\
\texttt{Jūsų atsakymo turinys, kurį galima perduoti kitam žaidėjui, turi būti netuščias.} \\
\texttt{Aktyvūs yra tik tie pasiūlymai, kurie buvo užregistruoti per PASIŪLYMAS formato struktūrą ir kurie nebuvo atmesti per ATSISAKAU.} \\
\\ 
            }
        }
    }
    & \\ \\

    \theutterance \stepcounter{utterance}  
    & \multicolumn{4}{p{0.6\linewidth}}{
        \cellcolor[rgb]{0.9,0.9,0.9}{
            \makecell[{{p{\linewidth}}}]{
                \texttt{\tiny{[P1$\rangle$GM]}}
                \texttt{PASIŪLYMAS: \{'A43', 'A67', 'B42', 'C10', 'B36', 'C61', 'C42'\}} \\
\\ 
\texttt{Bendros pastangos: 79 + 585 + 564 + 274 + 374 + 303 + 548 = 2727 (žemiau LIMIT = 3712)} \\
\texttt{Mano bendra svarba: 868 + 822 + 783 + 780 + 461 + 668 + 808 = 5190} \\
\\ 
\texttt{Šis rinkinys man yra labai vertingas ir palieka daug erdvės jūsų prioritetams. Kokie daiktai jums būtų svarbūs?} \\
            }
        }
    }
    & & \\ \\

    \theutterance \stepcounter{utterance}  
    & & \multicolumn{4}{p{0.6\linewidth}}{
        \cellcolor[rgb]{0.9,0.9,0.9}{
            \makecell[{{p{\linewidth}}}]{
                \texttt{\tiny{[P1$\langle$GM]}}
                \texttt{Naudojote žymę PASIŪLYMAS, bet nepateikėte galiojančios python aibės, kurioje kaip argumentai būtų eilutės, pvz., \{'A', 'B', 'C', ...\}. Bandykite dar kartą.} \\
            }
        }
    }
    & \\ \\

    \theutterance \stepcounter{utterance}  
    & \multicolumn{4}{p{0.6\linewidth}}{
        \cellcolor[rgb]{0.9,0.9,0.9}{
            \makecell[{{p{\linewidth}}}]{
                \texttt{\tiny{[P1$\rangle$GM]}}
                \texttt{PASIŪLYMAS: \{"A43", "A67", "B42", "C10", "B36", "C61", "C42"\}} \\
\\ 
\texttt{ARGUMENTAS: \{"Šis pasiūlymas palieka 985 pastangų jūsų prioritetams (3712 {-} 2727 = 985), o man užtikrina didelę naudą (5190 svarbos taškų)."\}} \\
\\ 
\texttt{Laukiu jūsų atsakymo ar kontrpasiūlymo.} \\
            }
        }
    }
    & & \\ \\

    \theutterance \stepcounter{utterance}  
    & & \multicolumn{4}{p{0.6\linewidth}}{
        \cellcolor[rgb]{0.9,0.9,0.9}{
            \makecell[{{p{\linewidth}}}]{
                \texttt{\tiny{[GM$\rangle$P2]}}
                \texttt{Dalyvaujate bendrų derybų žaidime.} \\
\\ 
\texttt{Kartu su kitu dalyviu turite susitarti dėl vieno daiktų rinkinio, kuris bus išsaugotas. Kiekvienas iš jūsų turi savo nuomonę, kiek kiekvienas daiktas jums svarbus (svarba). Jūs nežinote, kaip daiktus vertina kitas dalyvis. Be to, jums nurodoma, kiek pastangų reikalauja kiekvienas daiktas.} \\
\texttt{Dėl rinkinio galite susitarti tik tuo atveju, jei bendra pasirinktų daiktų pastangų suma neviršija bendros ribos:} \\
\\ 
\texttt{LIMIT = 3712} \\
\\ 
\texttt{Čia pateikiamos atskirų daiktų pastangų vertės:} \\
\\ 
\texttt{Daikto pastangos = \{"C08": 949, "C55": 388, "A43": 79, "A67": 585, "B42": 564, "A62": 230, "B19": 580, "C43": 84, "C73": 976, "C10": 274, "B36": 374, "A18": 913, "C61": 303, "C67": 578, "C42": 548\}} \\
\\ 
\texttt{Čia pateikiama jūsų asmeninė nuomonė apie kiekvieno daikto svarbą:} \\
\\ 
\texttt{Daikto svarbos reikšmės = \{"C08": 138, "C55": 583, "A43": 868, "A67": 822, "B42": 783, "A62": 65, "B19": 262, "C43": 121, "C73": 508, "C10": 780, "B36": 461, "A18": 484, "C61": 668, "C67": 389, "C42": 808\}} \\
\\ 
\texttt{Tikslas:} \\
\\ 
\texttt{Jūsų tikslas {-} susitarti dėl bendro daiktų rinkinio, kuris būtų jums kuo naudingesnis (t. y. maksimaliai padidintų bendrą svarbą JUMS), neviršijant LIMITO. Neprivalote kiekvienoje žinutėje pateikti PASIŪLYMO {-} galite tiesiog derėtis taip pat. Visos taktikos yra leidžiamos!} \\
\\ 
\texttt{Sąveikos protokolas:} \\
\\ 
\texttt{Žinutėje galite naudoti tik toliau nurodytus struktūrinius formatus:} \\
\\ 
\texttt{PASIŪLYMAS: \{'A', 'B', 'C', …\}} \\
\texttt{Siūlykite išlaikyti būtent šiuos daiktus.} \\
\texttt{ATSISAKAU: \{'A', 'B', 'C', …\}} \\
\texttt{Aiškiai atmesti oponento pasiūlymą.} \\
\texttt{ARGUMENTAS: \{'...'\}} \\
\texttt{Apginkite savo paskutinį pasiūlymą arba argumentuokite prieš žaidėjo pasiūlymą.} \\
\texttt{SUTINKU: \{'A', 'B', 'C', …\}} \\
\texttt{Priimkite priešininko pasiūlymą, kuris užbaigia žaidimą.} \\
\\ 
\\ 
\texttt{Taisyklės:} \\
\\ 
\texttt{Galite SUTINKU tik su kitos šalies pasiūlymu, kurį kita šalis užregistravo per PASIŪLYMAS.} \\
\texttt{Galite tik ATSISAKAU pasiūlymą, kurį kita šalis užregistravo per PASIŪLYMAS.} \\
\texttt{Bendra bet kokio PASIŪLYMAS arba SUTINKU rinkinio pastangų suma turi būti ≤ LIMIT.} \\
\texttt{NESKELBKITE savo paslėptų svarbos balų.} \\
\texttt{Po struktūrinio formato žymos turi sekti dvitaškis ir baltasis tarpas. Argumentas turi būti python rinkinys, kuriame yra 0 ar daugiau eilučių.} \\
\texttt{Taigi ji turi būti tokios formos: "TAG: \{...\}"} \\
\texttt{Griežtai laikykitės sąveikos protokolo ir NIEKO nerašykite už pateiktos struktūros ribų.} \\
\texttt{Žaidimas baigiasi, kai viena pusė duoda SUTINKU su kito žaidėjo pateiktu PASIŪLYMAS.} \\
\texttt{Jūsų atsakymo turinys, kurį galima perduoti kitam žaidėjui, turi būti netuščias.} \\
\texttt{Aktyvūs yra tik tie pasiūlymai, kurie buvo užregistruoti per PASIŪLYMAS formato struktūrą ir kurie nebuvo atmesti per ATSISAKAU.} \\
\\ 
\\ 
\\ 
\texttt{Dabar gausite pirmąjį kito žaidėjo pranešimą.} \\
\\ 
\texttt{PASIŪLYMAS: \{"A43", "A67", "B42", "C10", "B36", "C61", "C42"\}} \\
\\ 
\texttt{ARGUMENTAS: \{"Šis pasiūlymas palieka 985 pastangų jūsų prioritetams (3712 {-} 2727 = 985), o man užtikrina didelę naudą (5190 svarbos taškų)."\}} \\
\\ 
\texttt{Laukiu jūsų atsakymo ar kontrpasiūlymo.} \\
            }
        }
    }
    & \\ \\

    \theutterance \stepcounter{utterance}  
    & & & \multicolumn{4}{p{0.6\linewidth}}{
        \cellcolor[rgb]{0.9,0.9,0.9}{
            \makecell[{{p{\linewidth}}}]{
                \texttt{\tiny{[GM$\langle$P2]}}
                \texttt{ATSISAKAU: \{"A43", "A67", "B42", "C10", "B36", "C61", "C42"\}} \\
\\ 
\texttt{ARGUMENTAS: \{"Jūsų pasiūlymas man duoda tik 4468 svarbos taškų (A43:868 + A67:822 + B42:783 + C10:780 + B36:461 + C61:668 + C42:808), o aš turiu daiktų, kurie man yra svarbesni ir telpa į likusius 985 pastangų. Pavyzdžiui, C55 (583 svarba, 388 pastangos) arba C73 (508 svarba, 976 pastangos) būtų geresnis pasirinkimas nei kai kurie jūsų siūlomi daiktai. Siūlau peržiūrėti rinkinį, kad būtų įtraukti daugiau man svarbių daiktų."\}} \\
\\ 
\texttt{PASIŪLYMAS: \{"C55", "A43", "A67", "B42", "C10", "B36", "C61", "C43"\}} \\
            }
        }
    }
     \\ \\

    \theutterance \stepcounter{utterance}  
    & & \multicolumn{4}{p{0.6\linewidth}}{
        \cellcolor[rgb]{0.9,0.9,0.9}{
            \makecell[{{p{\linewidth}}}]{
                \texttt{\tiny{[P1$\langle$GM]}}
                \texttt{ATSISAKAU: \{"A43", "A67", "B42", "C10", "B36", "C61", "C42"\}} \\
\\ 
\texttt{ARGUMENTAS: \{"Jūsų pasiūlymas man duoda tik 4468 svarbos taškų (A43:868 + A67:822 + B42:783 + C10:780 + B36:461 + C61:668 + C42:808), o aš turiu daiktų, kurie man yra svarbesni ir telpa į likusius 985 pastangų. Pavyzdžiui, C55 (583 svarba, 388 pastangos) arba C73 (508 svarba, 976 pastangos) būtų geresnis pasirinkimas nei kai kurie jūsų siūlomi daiktai. Siūlau peržiūrėti rinkinį, kad būtų įtraukti daugiau man svarbių daiktų."\}} \\
\\ 
\texttt{PASIŪLYMAS: \{"C55", "A43", "A67", "B42", "C10", "B36", "C61", "C43"\}} \\
            }
        }
    }
    & \\ \\

    \theutterance \stepcounter{utterance}  
    & \multicolumn{4}{p{0.6\linewidth}}{
        \cellcolor[rgb]{0.9,0.9,0.9}{
            \makecell[{{p{\linewidth}}}]{
                \texttt{\tiny{[P1$\rangle$GM]}}
                \texttt{Jūsų pasiūlymas viršija LIMIT (bendros pastangos: 388 + 79 + 585 + 564 + 274 + 374 + 303 + 84 = **2651** + **C43** (84) jau įtrauktas, bet vis tiek per mažai vietos jums svarbiems daiktams, pvz., **C73** (976) netilptų).} \\
\\ 
\texttt{Siūlau tokią alternatyvą:} \\
\\ 
\texttt{**PASIŪLYMAS:** \{"C55", "A43", "A67", "B42", "C10", "B36", "C61"\}} \\
\texttt{Bendros pastangos: **2567** (lieka **1145** jums)} \\
\texttt{Mano svarba: **868 + 822 + 783 + 780 + 461 + 668 + 583 = 4965** (geresnis nei jūsų pateiktas 4468)} \\
\texttt{Jūsų naudai: Galite pridėti **C73** (976 pastangos) ir vis dar turėti **169** rezervą.} \\
\\ 
\texttt{**ARGUMENTAS:** \{"Šis rinkinys palieka daugiau erdvės jums įtraukti svarbius daiktus (pvz., C73) ir vis tiek garantuoja man aukštą naudą. Ar sutiktumėte?"\}} \\
            }
        }
    }
    & & \\ \\

    \theutterance \stepcounter{utterance}  
    & & & \multicolumn{2}{p{0.3\linewidth}}{
        \cellcolor[rgb]{0.9,0.9,0.9}{
            \makecell[{{p{\linewidth}}}]{
                \texttt{\tiny{[GM$|$GM]}}
                \texttt{abort} \\
            }
        }
    }
    & & \\ \\

    \theutterance \stepcounter{utterance}  
    & & & \multicolumn{2}{p{0.3\linewidth}}{
        \cellcolor[rgb]{0.9,0.9,0.9}{
            \makecell[{{p{\linewidth}}}]{
                \texttt{\tiny{[GM$|$GM]}}
                \texttt{end game} \\
            }
        }
    }
    & & \\ \\

\end{supertabular}
}

\end{document}
